\section{Introduction}

\emph{Deliveroo.js} is a grid-world game in which autonomous agents pick up
parcels and deliver them to delivery tiles for points. Parcel rewards decay
over time, actions are serialized by the server, and the world is only
partially observable through local sensing. The exam task is to build
\textbf{Agent~A}, a BDI agent for \emph{Challenge~1} (single-agent parcel
farming on eight maps), and \textbf{Agent~B}, an LLM-driven agent for
\emph{Challenge~2} (agents that shout natural-language missions), together with
Agent~A\,--\,Agent~B coordination and a meaningful \textbf{PDDL} planning
component.

We built a single ES-module codebase shared by both agents: belief revision,
option generation, swappable utility-based strategies, intention revision with
hysteresis, and a plan library with deterministic BFS movement on the directed
map graph, plus an optional PDDL plan for the \texttt{go\_to} sub-goal. Agent~B
adds, on top of the same machinery, an LLM mission interpreter and a structured
team protocol.

The headline Challenge~1 result is that a deliberately simple
\texttt{greedy-nearest} baseline is the strongest of four strategies across all
eight maps: these maps reward raw delivery throughput, and the value-aware
strategies lose because they deliver in small batches. We reached this
conclusion through a systematic benchmark and a log-level diagnosis that also
falsified our initial ``slow decay favours batching'' hypothesis --- a worked
example of using the metrics pipeline for analysis, not only scoring. The rest
of the report describes the architecture (Section~\ref{sec:arch}), Agent~A and
its strategies (Section~\ref{sec:bdi}), the LLM agent, PDDL and coordination
(Sections~\ref{sec:llm}--\ref{sec:coord}), and the experiments
(Section~\ref{sec:experiments}).
